\setcounter{section}{3}
\phantomsection
\section*{1. Introduction}
\addcontentsline{toc}{section}{1. Introduction}

There is a large body of literature that examines the effects of birth order. For economists in particular, research shows that whether you are the oldest or youngest or in between is not just a family curiosity but a factor with a distinct effect on education, skill development, and life success.

Empirical evidence using late twentieth- and early twenty-first century administrative data from Norway, the United States, and the United Kingdom consistently documents a firstborn premium, i.e., that firstborns are systematically better off than their later-born siblings with respect to school and other life outcomes \citep{Black2005,Behrman1986,Booth2009}. The reasons for this phenomenon are mainly explained by the allocation of time and money, as well as different expectations set for each child. That is why it is important to understand birth-order effects: they are the way we see families transferring advantage or disadvantage from one generation to the next.

However, evidence from historical populations in Europe suggests that birth-order effects may be limited or context-dependent. Using a dataset covering England between 1680 and 2024, \citet{ClarkCummins2024} find a small impact of birth order on social outcomes in the population over the long run.

Using Canadian survey data and U.S. administrative and educational data from the late twentieth century, recent studies comparing siblings within the same family show that younger children receive lower educational investments and attain lower levels of education than their older siblings \citep{Kantarevic2006, Hotz2015}.

Several mechanisms have been proposed to explain these patterns. One explanation is the dilution of parental resources in larger families. Another explanation is different parental investment and attention across siblings according to birth order. An additional explanation is sibling competition \citep{Becker1976, Blake1981, Price2008}. 

However, there is a gap in our historical understanding, since it is unclear whether these effects were also prevalent during pre-industrial times in Europe (from the thirteenth to the eighteenth centuries), as society was different from the contemporary one. In particular, families were larger, child mortality was higher, and legal norms, particularly those on inheritance, varied across regions. For these reasons, the impact of birth order in these historical contexts, and especially among elite groups such as scholars and intellectuals, remains an open question.

This dissertation focuses on this gap by looking at birth order effects amongst European scholars active in Europe during the aforementioned period. We combine biographical data on European scholars from the Repertorium Eruditorum Totius Europae (RETE), covering both universities and learned academies across Europe, with information on family relationships and regional inheritance systems. Focusing on scholars makes it possible to measure long-run academic success through publications and scholarly visibility instead of short-term educational outcomes.

One of the key components of this work is the role of institutional environments and inheritance systems. Inheritance institutions varied substantially across European regions between the sixteenth and eighteenth centuries \citep{Todd1990}. In Europe, there were primogeniture systems and egalitarian systems. Primogeniture systems concentrate family assets primarily on the eldest son, as opposed to egalitarian inheritance regimes that divide assets equally among siblings. These differences likely affect parental investment decisions and sibling competition. Recent research using historical data from France and other European regions between the eighteenth and twentieth centuries suggests a link between family systems (including inheritance institutions) and long-run demographic/socioeconomic outcomes \citep{Duranton2008, Gobbi2025}. However, it remains unclear to what extent these institutions generated inequalities within scholar families themselves.

Using new data on pairs of brothers who were both scholars, we pursue three main objectives. First, to determine if older brother scholars consistently outperformed their younger brothers and, if so, whether the size of this difference is similar to that found in modern studies. Second, to examine whether observed differences can be explained by factors such as age gaps and family background, and in particular whether birth-order productivity differences persist within shared institutional environments or instead become larger when siblings go to different institutions. Third, to examine how institutional environments and inheritance systems are associated with birth-order productivity differences within families. In particular, we hypothesize that the firstborn advantage will have been more pronounced in non-egalitarian systems.

The empirical strategy consists of within-family comparisons between brothers for whom birth order can be established from recorded birth years. This approach helps control for many characteristics shared within families, including parental background, social environment, and other family-level characteristics. Brother pairs with identical recorded birth years are excluded from the birth-order analysis, since older and younger siblings cannot be clearly identified in such cases. We measure scholarly productivity using publications and other historical indicators of academic visibility. We then examine how these outcomes relate to birth order and regional inheritance systems. The empirical approach combines paired comparisons with regression analysis.

This dissertation presents new evidence on the relationship between birth order and scholarly productivity among historical European scholars. We utilize a new dataset of 645 brother pairs to compare the performance of older and younger siblings, controlling for age gap, family background, and institutional context.

The findings provide evidence of a modest older-brother advantage in observed scholarly productivity among historical European scholars. However, this effect operates primarily through the extensive margin of academic activity. Older brothers are substantially more likely to have at least one recorded publication, whereas, conditional on publishing, productivity differences between siblings become small and statistically insignificant.

Using measures derived from VIAF, paired comparisons show that older brothers produced, on average, approximately two additional titles and tended to perform better across several indicators of academic prominence.

Last, we investigate whether inheritance systems and institutional environments are associated with different patterns of birth-order productivity gaps. The results suggest that institutional environments play a more robust role than inheritance systems in explaining productivity divergence between scholarly siblings. In particular, productivity differences become larger when brothers are affiliated with different academic institutions.

This dissertation contributes to several strands of the literature. First, it extends the literature on birth-order effects into the context of pre-industrial European elites by providing new evidence on scholarly productivity between the thirteenth and eighteenth centuries. Second, by combining RETE biographical data with information on family relationships, institutional environments, and inheritance systems, it contributes to the study of how historical academic environments shaped intra-family inequality in Europe. Third, the analysis highlights that birth-order differences among scholars operate mainly through entry into scholarly production rather than conditional productivity, thereby distinguishing between the extensive and intensive margins of intellectual activity.

The paper is organized as follows. Section~2 reviews the relevant literature on birth order, human capital formation, and family institutions. Section~3 describes the data, outlines the empirical strategy, and presents the main results together with their interpretation and limitations. Section~4 concludes.